\documentclass[11pt,a4paper]{article}

% ============================================================
%  Compilation : xelatex (deux passes)
% ============================================================
\usepackage{fontspec}
\setmainfont{Latin Modern Roman}
\usepackage[french]{babel}
\usepackage{amsmath, amsfonts, amssymb}
\usepackage{geometry}
\geometry{margin=2.3cm}
\usepackage{hyperref}
\usepackage{booktabs}
\usepackage{array}
\usepackage{tabularx}
\usepackage{enumitem}
\usepackage{float}
\usepackage[table]{xcolor}
\usepackage{caption}
\usepackage{tikz}
\usetikzlibrary{arrows.meta, positioning, fit, backgrounds, calc}

\definecolor{tabhead}{RGB}{31,78,121}
\definecolor{tabstripe}{RGB}{242,246,250}
\definecolor{tabchoice}{RGB}{222,238,212}
\definecolor{cReuse}{RGB}{214,230,245}
\definecolor{cNew}{RGB}{252,228,200}
\definecolor{cReal}{RGB}{225,225,225}
\definecolor{cOpt}{RGB}{235,222,245}
\newcolumntype{L}{>{\raggedright\arraybackslash}X}
\newcolumntype{P}[1]{>{\raggedright\arraybackslash}p{#1}}
\newcolumntype{B}[1]{>{\bfseries\raggedright\arraybackslash}p{#1}}
\newcommand{\hd}[1]{\textcolor{white}{\textbf{#1}}}
\newcommand{\choix}{\rowcolor{tabchoice}}
\captionsetup{font=small, labelfont={bf,color=tabhead}, labelsep=period}
\hypersetup{colorlinks=true, linkcolor=tabhead, urlcolor=tabhead}
\setlist{itemsep=2pt, topsep=4pt}

\newcommand{\decision}[1]{%
  \par\smallskip\noindent\colorbox{tabchoice}{\parbox{\dimexpr\linewidth-2\fboxsep}{\small\textbf{Décision à prendre --- }#1}}\par\smallskip}

\title{\textbf{Digital Twin Assisted Beamforming Design}\\[4pt]
\large Proposition d'architecture --- suite du projet SDM-VA-ISAC}
\author{}
\date{27 octobre 2026}

\begin{document}
\maketitle
\vspace{-1.2cm}
\tableofcontents
\newpage

% ============================================================
\section{Point de départ : où en est le projet}
% ============================================================

Le projet SDM-VA-ISAC a livré, sur la scène \texttt{phase3/new.xml} (quartier de Rabat, voiture
et deux piétons mobiles), une chaîne complète qui va du signal radio brut jusqu'à un jumeau
numérique sémantique :

\begin{itemize}
  \item \textbf{Phase 2} --- environnement Sionna RT 2.0.1 calibré (scène OSM, matériaux ITU, CIR,
        PDP, carte de couverture).
  \item \textbf{Phase 3} --- perception multimodale : nuage de points radar (4 n\oe uds BS + 2
        capteurs piétons), YOLOv8n fine-tuné sur les rendus (mAP50 = 0,965), fusion par
        \textit{cKDTree}.
  \item \textbf{Phase 4} --- grille bayésienne 3D, encodage hyperdimensionnel, visualisation
        Open3D/Plotly.
  \item \textbf{Phase 5} --- jumeau numérique ISAC double bande (FR1 3,5~GHz / 100~MHz, FR2
        28~GHz / 400~MHz, réseaux $8\times8$ et $16\times16$) : mesures radar \emph{apparentes}
        (fantômes multi-rebonds compris), filtrage par $\geq 2$ points de vue, objets créés par la
        vision. Configuration retenue C4 (FR1+FR2+vision) : F1 structures 0,52, voiture 100~\%,
        piéton A $77\pm13$~\%, piéton B $97\pm7$~\%, erreur de position des acteurs
        $\approx 0{,}35$--$0{,}52$~m, 0,08 faux objet par frame.
\end{itemize}

Dans la feuille de route d'origine, la phase suivante prévoyait un \og beamforming informé \fg{}
par la carte de risque. Le nouveau projet, \textbf{Digital Twin Assisted Beamforming Design}, en
fait le sujet principal. Le jumeau n'est plus une fin en soi : il devient \textbf{l'outil qui
calcule les faisceaux}.

\paragraph{Question scientifique proposée.}
\emph{Comment un jumeau numérique construit par perception ISAC et vision peut-il remplacer, en
partie ou en totalité, l'entraînement de faisceaux (\textit{beam training}) d'une BS mmWave, et
quelle fidélité du jumeau faut-il pour cela ?}

La seconde moitié de la question est celle qui relie les deux projets. La phase 5 a déjà produit
des jumeaux de fidélité différente (C1 FR1 seul, C2 FR2 seul, C3 double bande, C4 double bande +
vision). On peut donc mesurer directement \textbf{ce que chaque gain de perception rapporte en
qualité de faisceau}, une courbe que peu de travaux publient.

% ============================================================
\section{Principe général}
% ============================================================

Une BS FR2 doit orienter un faisceau étroit vers chaque utilisateur. Sans information a priori,
elle balaie un \textit{codebook} $\mathcal{W}=\{\mathbf{w}_1,\dots,\mathbf{w}_N\}$ ($N=256$ pour
un $16\times16$ en DFT) et garde le meilleur :
\begin{equation}
  \mathbf{w}^\star_t = \arg\max_{\mathbf{w}\in\mathcal{W}} \;\bigl|\mathbf{h}_t^{H}\mathbf{w}\bigr|^2 .
\end{equation}
Ce balayage coûte $N$ créneaux de référence (SSB/CSI-RS) et doit être refait dès que
l'utilisateur bouge ou qu'un obstacle coupe le trajet direct.

Le jumeau numérique sert de \textbf{simulateur de canal spécifique au site} :
\begin{enumerate}
  \item il connaît la géométrie statique, les objets mobiles (position, vitesse, classe) et la
        position des UE ;
  \item on reconstruit à partir de lui une \og scène jumelle \fg{} Sionna ;
  \item le ray tracing sur cette scène donne un canal prédit $\hat{\mathbf{h}}_t$, donc un
        classement des faisceaux \emph{avant} toute mesure ;
  \item la BS ne mesure plus que les $K \ll N$ meilleurs candidats.
\end{enumerate}
\begin{equation}
  \mathcal{S}_K = \operatorname{top-}K_{\mathbf{w}\in\mathcal{W}}\;
      \bigl|\hat{\mathbf{h}}_t^{H}\mathbf{w}\bigr|^2,
  \qquad
  \hat{\mathbf{w}}_t = \arg\max_{\mathbf{w}\in\mathcal{S}_K} \bigl|\mathbf{h}_t^{H}\mathbf{w}\bigr|^2 .
\end{equation}
Le compromis se lit sur le débit effectif, qui tient compte du temps perdu en entraînement :
\begin{equation}
  R_{\text{eff}} = \Bigl(1-\frac{K\,T_{\text{beam}}}{T_{\text{frame}}}\Bigr)
                   \log_2\!\Bigl(1+\frac{P\,|\mathbf{h}_t^{H}\hat{\mathbf{w}}_t|^2}{\sigma^2}\Bigr).
\end{equation}
Un jumeau parfait donne $K=1$ sans perte ; un jumeau faux oblige à augmenter $K$. Tout l'intérêt
est de quantifier la zone intermédiaire.

% ============================================================
\section{Architecture proposée}
% ============================================================

\begin{figure}[H]
\centering
\begin{tikzpicture}[
  font=\small,
  box/.style={draw=tabhead, thick, rounded corners=3pt, align=center,
              minimum height=1.25cm, text width=3.3cm, inner sep=4pt},
  reuse/.style={box, fill=cReuse},
  new/.style={box, fill=cNew},
  real/.style={box, fill=cReal},
  opt/.style={box, fill=cOpt, dashed},
  arr/.style={-{Stealth[length=2.4mm]}, thick, tabhead},
  darr/.style={-{Stealth[length=2.4mm]}, thick, violet!70!black, dashed},
  lab/.style={font=\scriptsize, fill=white, inner sep=1.5pt},
  node distance=1.2cm and 1.3cm]

  % Bande 1 : jumeau numérique
  \node[reuse] (dt) {\textbf{M1 --- État du jumeau}\\ phase 5 (C1--C4)\\ carte statique + $\mathcal{O}_t$};
  \node[new, right=of dt] (synth) {\textbf{M2 --- Scène jumelle}\\ bâtiments + gabarits\\ d'objets à $t+\Delta$};
  \node[new, right=of synth] (pred) {\textbf{M3 --- Prédicteur}\\ canal $\hat{\mathbf{h}}_{t+\Delta}$, gain\\ de chaque faisceau};

  % Bande 2 : beamforming
  \node[new, below=of dt] (track) {\textbf{M4 --- Prédiction}\\ trajectoires (Kalman)};
  \node[opt, below=of synth] (calib) {\textbf{M7 --- Calibration}\\ jumeau $\leftrightarrow$ réalité\\ (RT différentiable)};
  \node[new, below=of pred] (bf) {\textbf{M5 --- Conception}\\ choix de la BS\\ et du top-$K$ $\mathcal{S}_K$};

  % Bande 3 : réalité simulée
  \node[real, below=of track] (world) {\textbf{Monde \og réel \fg{}}\\ scène Sionna complète\\ (vérité terrain)};
  \node[real, below=of calib] (realch) {\textbf{Canal réel} $\mathbf{h}_t$\\ RT sur la scène vraie};
  \node[real, below=of bf] (sweep) {\textbf{M6 --- Mesure}\\ balayage réduit\\ sur $\mathcal{S}_K$ $\to$ $\hat{\mathbf{w}}_t$};

  % Flux principal
  \draw[arr] (world.west) -- ++(-0.75,0) |- node[lab, pos=0.25, rotate=90]{capteurs ISAC + caméras} (dt.west);
  \draw[arr] (dt) -- (synth);
  \draw[arr] (synth) -- (pred);
  \draw[arr] (dt) -- (track);
  \draw[arr] (track.north east) -- node[lab]{$\hat{\mathbf{p}}_{t+\Delta}$} (synth.south west);
  \draw[arr] (pred) -- node[lab]{gains prédits} (bf);
  \draw[arr] (bf) -- node[lab]{$\mathcal{S}_K$} (sweep);
  \draw[arr] (world) -- (realch);
  \draw[arr] (realch) -- (sweep);

  % Boucle optionnelle
  \draw[darr] (sweep.north west) -- node[lab]{écart} (calib.south east);
  \draw[darr] (calib) -- node[lab]{matériaux} (synth);

  \begin{scope}[on background layer]
    \node[fit=(dt)(pred), draw=tabhead!60, rounded corners=6pt, inner sep=6pt, fill=cReuse!25,
          label={[font=\scriptsize\bfseries, tabhead, anchor=west, align=left]east:Jumeau\\numérique}] {};
    \node[fit=(track)(bf), draw=orange!70!black, rounded corners=6pt, inner sep=6pt, fill=cNew!25,
          label={[font=\scriptsize\bfseries, orange!70!black, anchor=west, align=left]east:Beam-\\forming}] {};
    \node[fit=(world)(sweep), draw=gray, rounded corners=6pt, inner sep=6pt, fill=gray!6,
          label={[font=\scriptsize\bfseries, gray, anchor=west, align=left]east:Réalité\\simulée}] {};
  \end{scope}
\end{tikzpicture}
\caption{Architecture proposée. Bleu : réutilisé des phases 3--5. Orange : nouveau. Violet en
pointillés : boucle optionnelle de calibration, alimentée par l'écart entre gains prédits et
mesurés. Gris : la \og réalité \fg{}, simulée par une scène Sionna distincte
de la scène jumelle.}
\label{fig:archi}
\end{figure}

\subsection{Séparer la réalité du jumeau}

C'est le point méthodologique le plus important. Tout est simulé avec Sionna, donc si la scène
jumelle était identique à la scène \og réelle \fg{}, le jumeau serait parfait par construction et
les résultats n'auraient aucune valeur. On maintient donc \textbf{deux scènes distinctes} :

\begin{itemize}
  \item \textbf{Scène réelle} : \texttt{new.xml} complète, acteurs à leur vraie position, diffusion
        diffuse activée, matériaux éventuellement perturbés (écart sur $\varepsilon_r$ et
        $\sigma$ autour des valeurs ITU), petits objets absents du jumeau (lampadaires, mobilier).
  \item \textbf{Scène jumelle} : construite \emph{uniquement} à partir de ce que le jumeau sait
        (sortie de la phase 5) --- positions estimées avec leur erreur, classes reconnues,
        objets manqués absents, faux objets présents.
\end{itemize}
L'écart entre les deux (\textit{twin-reality gap}) est la variable que l'on fait varier.

\subsection{Rôle de chaque module}

\begin{table}[H]
\centering
\caption{Modules de l'architecture.}
\label{tab:modules}
\rowcolors{2}{white}{tabstripe}
\begin{tabularx}{\linewidth}{B{2.6cm} L P{2.2cm}}
\rowcolor{tabhead}
\hd{Module} & \hd{Rôle} & \hd{Origine} \\
M1 État du jumeau & Carte statique (voxels/classe), objets $\mathcal{O}_t=\{(\mathbf{p}_i,\mathbf{v}_i,c_i,\boldsymbol{\Sigma}_i)\}$, positions UE estimées. Plusieurs niveaux de fidélité (C1--C4). & Phase 5, inchangé \\
M2 Scène jumelle & Transforme l'état en scène Sionna : bâtiments (maillage OSM ou voxels) + gabarit par classe (boîte $4{,}5\times1{,}8\times1{,}5$~m pour une voiture, cylindre pour un piéton) placé à la position estimée, matériau ITU selon la classe. & Nouveau \\
M3 Prédicteur de canal & Ray tracing sur la scène jumelle $\to$ $\hat{\mathbf{h}}$ par couple BS--UE, puis gain prédit de chaque faisceau du codebook. & Nouveau \\
M4 Prédiction & Filtre de Kalman sur les trajectoires (vitesse déjà estimée en FR2 par Doppler) $\to$ positions à $t+\Delta$, qui permettent de repérer un blocage du trajet direct avant qu'il ne se produise. & Nouveau \\
M5 Conception des faisceaux & Choix de la BS (4 disponibles) et du sous-ensemble $\mathcal{S}_K$ ; commutation anticipée avant un blocage. & Nouveau \\
M6 Mesure & Canal vrai sur la scène réelle, balayage limité à $\mathcal{S}_K$, calcul de la perte et du débit effectif. & Nouveau \\
M7 Calibration & Ajuste les matériaux (et décalages de position) de la scène jumelle par gradient, en comparant puissances prédites et mesurées (Sionna RT est différentiable). & Optionnel \\
\end{tabularx}
\end{table}

% ============================================================
\section{Choix d'architecture}
% ============================================================

Quatre décisions structurent le projet. Pour chacune, je donne les options et ma recommandation ;
c'est à toi de trancher.

\subsection{D1 --- Paradigme de beamforming (décision principale)}

\begin{table}[H]
\centering
\caption{Trois paradigmes possibles. La ligne verte est la recommandation.}
\label{tab:d1}
\rowcolors{2}{white}{tabstripe}
\begin{tabularx}{\linewidth}{B{3.0cm} L L P{1.5cm}}
\rowcolor{tabhead}
\hd{Option} & \hd{Principe} & \hd{Intérêt / risque} & \hd{Charge} \\
\choix A. Sélection de faisceau assistée par jumeau (+ prédictive) &
Codebook analogique DFT, un faisceau par UE et par BS. Le jumeau classe les faisceaux, la BS ne
mesure que le top-$K$. Extension : prédiction à $t+\Delta$ et commutation de BS avant blocage. &
Continuité directe avec la phase 5 (positions + objets), métriques standard et lisibles, résultat
garanti. Moins \og algorithmique \fg{} que B ou C. & Moyenne \\
B. Précodage hybride multi-utilisateur &
Le jumeau donne $\hat{\mathbf{H}}$ complet ; étage analogique par sélection (A), étage numérique
ZF/MMSE/WMMSE sur le canal effectif, plusieurs UE servis simultanément. &
Plus proche d'une vraie conception de précodeur ; sensible à la phase du canal prédit, que le
jumeau estime mal (erreur de 0,4~m $\gg \lambda = 1{,}07$~cm). & Élevée \\
C. Beam management appris (DRL) dans le jumeau &
Un agent apprend la politique (BS, faisceau, moment du balayage) avec un état enrichi par la
carte sémantique $\boldsymbol{\xi}_t^{\text{SDM}}$ ; entraînement dans le jumeau, test sur la
réalité. &
Lien fort avec le \textit{Telecom Brain} et la feuille de route d'origine. Demande beaucoup de
trajectoires, donc beaucoup de ray tracing ; risque de ne pas converger dans les délais. & Très
élevée \\
\end{tabularx}
\end{table}

\paragraph{Recommandation : A comme c\oe ur, B en extension.}
L'argument décisif est physique. À 28~GHz, une erreur de position de 0,4~m vaut environ 37
longueurs d'onde : le jumeau ne peut pas prédire la \emph{phase} du canal, seulement sa
\emph{structure} (directions des trajets, puissances, blocages). La sélection de faisceau n'a
besoin que de cette structure ; le précodage numérique (B) a besoin de la phase et devra de toute
façon s'appuyer sur une mesure. B s'ajoute donc naturellement \emph{par-dessus} A : le jumeau
choisit les faisceaux analogiques, une mesure courte du canal effectif sert au précodeur
numérique. C reste possible en fin de projet si le temps le permet, en réutilisant M1--M6 comme
environnement d'entraînement.

\decision{A seul, A puis B, ou C ?}

\subsection{D2 --- Construction de la scène jumelle}

\begin{itemize}
  \item \textbf{D2-a. Carte statique connue + objets perçus} \emph{(recommandé)}. Les bâtiments
        viennent du modèle 3D (OSM, comme le ferait un opérateur), seuls les objets mobiles
        viennent de la perception. Réaliste : un opérateur dispose d'un modèle de ville, ce qu'il
        ne connaît pas, ce sont les voitures et les piétons.
  \item \textbf{D2-b. Tout depuis la perception}. Les bâtiments sont reconstruits à partir des
        voxels de la phase 5 (marching cubes ou boîtes par voxel). Plus pur scientifiquement, mais
        le rappel structures de 0,36 donnerait des façades trouées et des trajets réfléchis
        absents, ce qui pénaliserait surtout la reconstruction et non le beamforming.
\end{itemize}
On peut faire D2-a comme référence et D2-b comme expérience complémentaire, ce qui ajoute un
niveau à l'échelle de fidélité ci-dessous.

\decision{D2-a seul, ou D2-a + D2-b ?}

\subsection{D3 --- Calcul du canal prédit}

\begin{itemize}
  \item \textbf{D3-a. Ray tracing en ligne à chaque frame} : exact, mais un appel
        \texttt{PathSolver} par frame et par BS.
  \item \textbf{D3-b. Carte de canal précalculée + correction de blocage} \emph{(recommandé)} :
        on précalcule une fois sur la scène statique une carte de connaissance du canal
        (\textit{channel knowledge map}) sur une grille de positions UE (meilleur faisceau et gain
        par faisceau, par BS) ; en ligne, on ne fait qu'un test de visibilité contre les gabarits
        des objets mobiles et on retire les trajets bloqués. Rapide, et c'est l'architecture
        qu'un opérateur déploierait réellement.
\end{itemize}
D3-a sert de référence pour mesurer ce que l'approximation D3-b coûte en précision.

\decision{D3-b avec D3-a comme référence, ou D3-a uniquement ?}

\subsection{D4 --- Bandes de fréquence}

Recommandation : \textbf{jumeau construit en FR1+FR2+vision (C4), beamforming en FR2}. C'est en
FR2 que le beamforming est critique (faisceaux étroits, blocage par un piéton) et que
l'entraînement coûte cher ($N=256$). FR1 continue de servir au sensing, comme en phase 5.
Alternative : beamforming aussi en FR1 ($8\times8$, $N=64$) pour montrer que le bénéfice du jumeau
croît avec la fréquence.

\decision{FR2 seul, ou FR1 et FR2 en comparaison ?}

% ============================================================
\section{Scénario de simulation}
% ============================================================

\begin{itemize}
  \item \textbf{Scène} : \texttt{phase3/new.xml} (déjà validée, matériaux corrigés pour 28~GHz).
  \item \textbf{BS} : les 4 n\oe uds TX1--TX4 existants, réseau UPA $16\times16$ à 28~GHz,
        codebook DFT $N=256$ (éventuellement suréchantillonné).
  \item \textbf{UE} : le piéton A, le piéton B et la voiture portent chacun un UE (les acteurs sont
        à la fois utilisateurs et obstacles), plus 2--4 UE piétons supplémentaires sur les
        trottoirs pour multiplier les situations de blocage.
  \item \textbf{Séquence} : la phase 5 n'a que 6 frames. Il faut une trajectoire plus longue, par
        exemple 10~s à 10~frames/s (100 frames), pour que la prédiction et les blocages aient un
        sens. C'est le principal coût de calcul nouveau.
  \item \textbf{Jumeaux évalués} : on fait varier la fidélité selon l'échelle suivante.
\end{itemize}

\begin{table}[H]
\centering
\caption{Échelle de fidélité du jumeau : la variable principale des expériences.}
\label{tab:fidelite}
\rowcolors{2}{white}{tabstripe}
\begin{tabularx}{\linewidth}{B{1.2cm} P{4.4cm} L}
\rowcolor{tabhead}
\hd{Niveau} & \hd{Contenu de la scène jumelle} & \hd{Ce qu'il mesure} \\
J0 & Aucune scène : faisceau géométrique vers la position estimée de l'UE (ligne de vue supposée) & Ce que donne la seule localisation, sans jumeau \\
J1 & Bâtiments seuls (OSM) & Apport des réflexions sur les façades \\
J2 & J1 + objets perçus, configurations C1 à C4 de la phase 5 & \textbf{Apport de chaque gain de perception sur le faisceau} \\
J3 & J2 (C4) + matériaux calibrés (M7) & Apport de la boucle de calibration \\
J$\infty$ & Scène réelle exacte (oracle) & Borne supérieure atteignable par un jumeau \\
\end{tabularx}
\end{table}

% ============================================================
\section{Évaluation}
% ============================================================

\subsection{Références de comparaison}
\begin{itemize}
  \item \textbf{Balayage exhaustif} des $N$ faisceaux : performance optimale, coût maximal.
  \item \textbf{Recherche hiérarchique} (faisceaux larges puis fins, $\approx 2\log_2 N$ mesures) :
        référence standard 5G NR.
  \item \textbf{J0, faisceau par localisation seule} : montre ce que la carte apporte au-delà de
        la position.
  \item \textbf{Suivi réactif} : on ne rebalaie qu'après une chute de puissance (pas d'anticipation
        des blocages).
\end{itemize}

\subsection{Métriques}
\begin{table}[H]
\centering
\caption{Métriques proposées.}
\label{tab:metriques}
\rowcolors{2}{white}{tabstripe}
\begin{tabularx}{\linewidth}{B{4.2cm} L}
\rowcolor{tabhead}
\hd{Métrique} & \hd{Définition} \\
Précision top-$K$ & Fraction des frames où le meilleur faisceau réel appartient à $\mathcal{S}_K$ \\
Perte de SNR & $10\log_{10}\bigl(|\mathbf{h}^H\mathbf{w}^\star|^2 / |\mathbf{h}^H\hat{\mathbf{w}}|^2\bigr)$, médiane et 90\textsuperscript{e} centile \\
Surcoût d'entraînement & $K/N$, ou $K$ nécessaire pour une perte $< 1$~dB dans 90~\% des cas \\
Débit effectif & $R_{\text{eff}}$ (équation~3), moyenné sur la séquence \\
Coupures & Nombre et durée des frames sous un seuil de SNR, surtout lors des blocages \\
Anticipation & Délai entre la commutation décidée par le jumeau et le blocage réel \\
Coût de calcul & Temps par frame de M2--M5 sur la RTX 4070, comparé à la durée d'une frame \\
\end{tabularx}
\end{table}

\paragraph{Résultat attendu principal.} Une courbe \og fidélité du jumeau $\to$ surcoût
d'entraînement \fg{} : pour chaque niveau J0 à J$\infty$, le $K$ minimal qui garantit une perte de
SNR inférieure à 1~dB. Elle relie directement les résultats de la phase 5 (F1, erreur de position,
taux de reconnaissance des piétons) au gain système, ce qui donne au jumeau une valeur mesurable
en termes de réseau.

% ============================================================
\section{Plan de travail proposé}
% ============================================================

\begin{table}[H]
\centering
\caption{Étapes, dans l'ordre de dépendance.}
\label{tab:plan}
\rowcolors{2}{white}{tabstripe}
\begin{tabularx}{\linewidth}{B{0.6cm} P{4.2cm} L}
\rowcolor{tabhead}
\hd{\#} & \hd{Étape} & \hd{Livrable} \\
1 & Séquence longue + UE & Trajectoires des acteurs et UE sur 100 frames, sortie du jumeau phase 5 pour chaque frame \\
2 & Canal réel et codebook & M6 : $\mathbf{h}_t$ sur la scène réelle pour chaque BS--UE, gains des $N$ faisceaux, faisceau optimal (vérité terrain) \\
3 & Scène jumelle & M2 : génération automatique de la scène Sionna à partir de l'état du jumeau (D2) \\
4 & Prédicteur et sélection & M3 + M5 statiques : précision top-$K$ et perte de SNR pour J0, J1, J2, J$\infty$ \\
5 & Prédiction et blocages & M4 : Kalman, commutation anticipée de BS/faisceau, comparaison au suivi réactif \\
6 & Extension (selon D1) & B : précodage numérique multi-UE ; ou M7 : calibration des matériaux (J3) \\
7 & Évaluation et rapport & Courbe fidélité $\to$ surcoût, tableau des références, discussion \\
\end{tabularx}
\end{table}

% ============================================================
\section{Limites à assumer dès le départ}
% ============================================================
\begin{itemize}
  \item \textbf{Réalité simulée.} La \og réalité \fg{} est elle aussi un ray tracing Sionna. On
        contrôle l'écart jumeau/réalité, mais on ne capture pas tout ce qu'un vrai canal
        contient. Les perturbations de matériaux et les objets absents du jumeau en atténuent
        l'effet, sans le supprimer.
  \item \textbf{Phase du canal.} Le jumeau ne prédira pas la phase à 28~GHz (cf. D1) ; c'est
        une limite physique, pas d'implémentation, et elle justifie de garder un court balayage.
  \item \textbf{Variabilité GPU.} La diffusion diffuse varie d'un run à l'autre, comme en phase 5 ;
        on garde la même pratique (plusieurs tirages, moyenne $\pm$ écart-type).
  \item \textbf{TX4} n'a pas de vue caméra sur les acteurs : ce n'est pas bloquant ici, la position
        des objets venant de la fusion multi-n\oe uds, mais TX4 reste un candidat de commutation
        intéressant justement parce que sa géométrie est différente.
\end{itemize}

\vspace{0.4cm}
\noindent\colorbox{tabchoice}{\parbox{\dimexpr\linewidth-2\fboxsep}{%
\textbf{Synthèse des recommandations.} D1 : sélection de faisceau assistée par jumeau, prédictive
et sensible aux blocages (A), avec précodage hybride (B) en extension. D2 : carte statique connue
+ objets perçus. D3 : carte de canal précalculée + correction de blocage, avec ray tracing en
ligne comme référence. D4 : jumeau FR1+FR2+vision, beamforming en FR2.}}

\end{document}
